% Lean module: MonadFormal.Temporal.Variational
\subsection{A coercive whole-path energy}
\label{sec:temporal-variational}

Whole-path consistency is a third interpretation of temporal dependence.
It specifies an extremum subject to boundary data, rather than an
initial-value update rule. The following Hilbert-space result supplies
the existence argument for the interval construction below.

Let $V$ be a real Hilbert space, $W$ a real inner-product space,
$T:V\to W$ a bounded linear operator, and $b\in W$. Put
\[
 E_\lambda(u)=\frac12\lVert u\rVert^2
       +\frac\lambda2\lVert Tu+b\rVert^2,\qquad
 B_\lambda(u,h)=\langle u,h\rangle+\lambda\langle Tu,Th\rangle.
\]
\leandef{MonadFormal.Temporal.Variational.energy,MonadFormal.Temporal.Variational.form}

\begin{theorem}[Coercive quadratic minimization]
\label{thm:variational-minimum}
If $\lambda\ge0$, then $E_\lambda$ has a unique global minimizer.
There is a stationary point $u_*$ satisfying
\begin{equation}\label{eq:variational-weak}
 \langle u_*,h\rangle+\lambda\langle Tu_*+b,Th\rangle=0
 \quad\text{for every }h\in V.
\end{equation}
Every such point obeys the exact gap identity
\begin{equation}\label{eq:variational-gap}
 E_\lambda(u_*+h)-E_\lambda(u_*)
 =\frac12\bigl(\lVert h\rVert^2+\lambda\lVert Th\rVert^2\bigr).
\end{equation}
\end{theorem}
\begin{proof}
The bounded bilinear form satisfies
$B_\lambda(u,u)\ge\lVert u\rVert^2$.
The functional $h\mapsto-\lambda\langle b,Th\rangle$ is bounded,
so Lax--Milgram gives \eqref{eq:variational-weak}.
Expand the two squared norms at $u_*+h$; the linear cross terms
vanish by stationarity, giving \eqref{eq:variational-gap}.
Its right-hand side is nonnegative and is zero only if $h=0$.
Thus $u_*$ minimizes the energy, and any other minimizer equals it.
\leanproof{MonadFormal.Temporal.Variational.form_apply,MonadFormal.Temporal.Variational.form_coercive,MonadFormal.Temporal.Variational.exists_stationary,MonadFormal.Temporal.Variational.energy_gap,MonadFormal.Temporal.Variational.minimizes_of_stationary,MonadFormal.Temporal.Variational.exists_unique_minimizer}
\end{proof}

% Lean module: MonadFormal.Temporal.Variational.Interval
\subsubsection{Endpoint-constrained interval paths}
\label{sec:variational-interval}

Let $\mu$ be Lebesgue measure restricted to $[0,1]$, and set
$V=L^2(\mu;\mathbb R)$. For $f\in V$, define
\[
 (Pf)(t)=\int_0^t f(s)\,d\mu(s),\qquad
 H=\{f\in V:(Pf)(1)=0\}.
\]
Functions in $V$ are identified almost everywhere. Their primitives
are ordinary continuous functions; the integral does not depend on
the representative. Given endpoint values $a,b\in\mathbb R$, put
\[
 x_f(t)=a+(b-a)t+(Pf)(t),\qquad f\in H.
\]
The restricted measure extends the primitive constantly outside the
unit interval; this is an integration convention, not an additional
equation of motion.
\leandef{MonadFormal.Temporal.Variational.Interval.unitMeasure,MonadFormal.Temporal.Variational.Interval.Velocity,MonadFormal.Temporal.Variational.Interval.primitive,MonadFormal.Temporal.Variational.Interval.primitiveCLM,MonadFormal.Temporal.Variational.Interval.endpoint,MonadFormal.Temporal.Variational.Interval.variationSpace,MonadFormal.Temporal.Variational.Interval.path}

\begin{lemma}[Bounded primitives and endpoint variations]
\label[lemma]{lem:variational-paths}
The map $P:V\to C_b(\mathbb R)$ is bounded linear, with
$\lVert Pf\rVert_\infty\le\lVert f\rVert_2$.
The subspace $H$ is closed. Each $x_f$ is absolutely continuous on
$[0,1]$, has endpoints $x_f(0)=a$, $x_f(1)=b$, and satisfies
$x_f'=b-a+f$ almost everywhere on that interval. Moreover,
\[
 \int f\,d\mu=0,\qquad
 \int (b-a+f)^2\,d\mu=(b-a)^2+\lVert f\rVert_2^2
 \quad(f\in H).
\]
\end{lemma}
\begin{proof}
Since $\mu(\mathbb R)=1$, Cauchy--Schwarz bounds each oriented
primitive by $\lVert f\rVert_2$. Integration is linear, and an
integrable function has a continuous primitive. Evaluation at $1$
is bounded, so its composite with $P$ has closed kernel $H$.
Extend $f$ by zero outside $[0,1]$. Ordinary integration of this
extension equals $Pf$. The fundamental theorem for Lebesgue
integrals gives absolute continuity and derivative $f$ almost
everywhere. Add the affine path to obtain the endpoint and derivative
claims. Finally, $(Pf)(1)=\int f\,d\mu=0$ removes the cross term
in the expansion of $(b-a+f)^2$.
\leanproof{MonadFormal.Temporal.Variational.Interval.velocity_integrable,MonadFormal.Temporal.Variational.Interval.primitive_continuous,MonadFormal.Temporal.Variational.Interval.primitive_bound,MonadFormal.Temporal.Variational.Interval.primitive_add,MonadFormal.Temporal.Variational.Interval.primitive_smul,MonadFormal.Temporal.Variational.Interval.primitiveCLM_apply,MonadFormal.Temporal.Variational.Interval.variation_initial,MonadFormal.Temporal.Variational.Interval.variation_terminal,MonadFormal.Temporal.Variational.Interval.extendedVelocity_integrable,MonadFormal.Temporal.Variational.Interval.primitive_eq_integral,MonadFormal.Temporal.Variational.Interval.primitive_absolutelyContinuous,MonadFormal.Temporal.Variational.Interval.primitive_hasDerivAt_ae,MonadFormal.Temporal.Variational.Interval.path_endpoints,MonadFormal.Temporal.Variational.Interval.path_hasDerivAt_ae,MonadFormal.Temporal.Variational.Interval.path_absolutelyContinuous,MonadFormal.Temporal.Variational.Interval.primitive_one_eq_integral,MonadFormal.Temporal.Variational.Interval.variation_integral_zero,MonadFormal.Temporal.Variational.Interval.integral_sq_eq_norm_sq,MonadFormal.Temporal.Variational.Interval.kinetic_integral}
\end{proof}

Fix $0\le\delta\le1$ and let $\nu_\delta$ be Lebesgue measure
restricted to $[0,1-\delta]$. Define
\[
 T_\delta f(t)=Pf(t+\delta)-Pf(t),\qquad
 b_\delta(t)=(b-a)\delta.
\]
Here $T_\delta:H\to L^2(\nu_\delta;\mathbb R)$ is bounded:
it composes $P$, translation minus the identity, and inclusion of
bounded continuous functions into $L^2$ of a finite measure.
\leandef{MonadFormal.Temporal.Variational.Interval.shiftMeasure,MonadFormal.Temporal.Variational.Interval.differenceCLM,MonadFormal.Temporal.Variational.Interval.shiftCLM,MonadFormal.Temporal.Variational.Interval.affineOffset,MonadFormal.Temporal.Variational.Interval.intervalEnergy}

\begin{theorem}[Nonlocal minimization in velocity coordinates]
\label{thm:interval-minimum}
For $\lambda\ge0$, the energy
\[
 J(x_f)=\frac12\int (b-a+f)^2\,d\mu
  +\frac\lambda2\int\bigl(x_f(t+\delta)-x_f(t)\bigr)^2\,d\nu_\delta(t)
\]
has a unique minimizing velocity class $f_*\in H$. It satisfies
\[
 \langle f_*,h\rangle_{L^2(\mu)}
 +\lambda\langle T_\delta f_*+b_\delta,T_\delta h\rangle_{L^2(\nu_\delta)}=0
 \quad(h\in H).
\]
\end{theorem}
\begin{proof}
Almost everywhere,
$T_\delta f+b_\delta=x_f(\,\cdot+\delta)-x_f$.
The preceding lemma therefore gives
\[
 J(x_f)=\frac{(b-a)^2}{2}
       +\frac{\lVert f\rVert_2^2}{2}
       +\frac\lambda2\lVert T_\delta f+b_\delta\rVert_2^2.
\]
Apply \Cref{thm:variational-minimum} to the Hilbert space $H$,
operator $T_\delta$ and offset $b_\delta$. The constant term does
not affect minimization.
\leanproof{MonadFormal.Temporal.Variational.Interval.shiftCLM_apply_ae,MonadFormal.Temporal.Variational.Interval.shifted_path_ae,MonadFormal.Temporal.Variational.Interval.intervalEnergy_eq_integrals,MonadFormal.Temporal.Variational.Interval.exists_unique_interval_minimizer,MonadFormal.Temporal.Variational.Interval.exists_interval_stationary_minimizer}
\end{proof}

% Lean module: MonadFormal.Temporal.Variational.Paths
\subsubsection{The full admissible path class}
\label{sec:variational-admissible}

Let $\mathcal A_{a,b}$ consist of absolutely continuous functions
$x:[0,1]\to\mathbb R$ with $x'\in L^2(\mu)$ and endpoints
$x(0)=a$, $x(1)=b$. These are the absolutely continuous
representatives of the endpoint-constrained one-dimensional $H^1$
class. Define
\[
 J(x)=\frac12\int (x')^2\,d\mu
      +\frac\lambda2\int\bigl(x(t+\delta)-x(t)\bigr)^2\,d\nu_\delta(t).
\]
Only values on $[0,1]$ enter this functional. In particular, extending
a path outside that interval changes neither its derivative almost
everywhere on the interval nor its shifted difference on
$[0,1-\delta]$.
\leandef{MonadFormal.Temporal.Variational.Interval.Admissible,MonadFormal.Temporal.Variational.Interval.pathEnergy}

\begin{theorem}[Exact velocity coordinates]
\label{thm:variational-coordinates}
The map $f\mapsto x_f|_{[0,1]}$ is a bijection from $H$ onto
$\mathcal A_{a,b}$. Its inverse is the $L^2$ class
$f=x'-(b-a)$. The path energy equals the coordinate energy in
\Cref{thm:interval-minimum}.
\end{theorem}
\begin{proof}
The construction gives an admissible path by
\Cref{lem:variational-paths}. Conversely, if $x\in\mathcal A_{a,b}$,
then $f=x'-(b-a)$ belongs to $L^2(\mu)$. Absolute continuity and
the fundamental theorem of calculus give, for $0\le t\le1$,
\[
 (Pf)(t)=x(t)-a-(b-a)t.
\]
At $t=1$ this is zero, so $f\in H$ and $x=x_f$ on $[0,1]$.
If two constructed paths agree there, their derivatives agree at
every interior point where both derivatives exist. Their velocity
classes therefore agree almost everywhere, giving uniqueness.
Finally, $x_f'=b-a+f$ almost everywhere identifies the kinetic
integral, while the shifted-path integral is unchanged.
\leanproof{MonadFormal.Temporal.Variational.Interval.unit_integral_eq,MonadFormal.Temporal.Variational.Interval.ae_mem_interior,MonadFormal.Temporal.Variational.Interval.deriv_eq_ae_of_eqOn,MonadFormal.Temporal.Variational.Interval.path_injective,MonadFormal.Temporal.Variational.Interval.path_admissible,MonadFormal.Temporal.Variational.Interval.primitive_centered_derivative,MonadFormal.Temporal.Variational.Interval.exists_unique_velocity,MonadFormal.Temporal.Variational.Interval.pathEnergy_path,MonadFormal.Temporal.Variational.Interval.pathEnergy_congr}
\end{proof}

\begin{theorem}[Unique admissible minimizing path]
\label{thm:variational-path-minimum}
For $0\le\delta\le1$ and $\lambda\ge0$, $J$ has a unique
minimizer in $\mathcal A_{a,b}$.
\end{theorem}
\begin{proof}
Transfer the minimizing velocity from \Cref{thm:interval-minimum}
through the energy-preserving bijection of
\Cref{thm:variational-coordinates}. Every admissible competitor
has a velocity coordinate, and every minimizing path must have
the unique minimizing coordinate. Thus uniqueness is equality
on $[0,1]$, with no assertion about an extension outside it.
\leanproof{MonadFormal.Temporal.Variational.Interval.exists_unique_path_minimizer}
\end{proof}

% Lean module: MonadFormal.Temporal.Variational.Shift
\subsubsection{The shifted-difference adjoint}
\label{sec:variational-shift}

For $0\le\delta\le1$, define the nonlocal force
\begin{equation}\label{eq:variational-force}
 G_\delta[x](t)
 =\mathbf1_{[\delta,1]}(t)\bigl(x(t)-x(t-\delta)\bigr)
  -\mathbf1_{[0,1-\delta]}(t)\bigl(x(t+\delta)-x(t)\bigr).
\end{equation}
The indicators are part of the operator. They exclude a shifted
neighbour that lies outside the interval; no periodic or exterior
boundary value is supplied.
\leandef{MonadFormal.Temporal.Variational.Interval.shiftForce}

\begin{lemma}[Shifted pairing with its boundary supports]
\label[lemma]{lem:shift-pairing}
For continuous real functions $x,h$,
$G_\delta[x]\in L^2(\mu)$ and
\[
 \int G_\delta[x](t)h(t)\,d\mu(t)
 =\int\bigl(x(t+\delta)-x(t)\bigr)
          \bigl(h(t+\delta)-h(t)\bigr)\,d\nu_\delta(t).
\]
\end{lemma}
\begin{proof}
Continuous functions are bounded on $[0,1]$, so the two restricted
differences are square-integrable. Split the product on the right.
In the term containing $h(t+\delta)$, set $s=t+\delta$:
\[
 \int_0^{1-\delta}\!\bigl(x(t+\delta)-x(t)\bigr)h(t+\delta)\,dt
 =\int_\delta^1\!\bigl(x(s)-x(s-\delta)\bigr)h(s)\,ds.
\]
The other term remains on $[0,1-\delta]$. Writing both as
integrals on $[0,1]$ introduces precisely the two indicators in
\eqref{eq:variational-force}. Endpoints have Lebesgue measure zero.
\leanproof{MonadFormal.Temporal.Variational.Interval.continuous_memLp_unit,MonadFormal.Temporal.Variational.Interval.shiftForce_memLp,MonadFormal.Temporal.Variational.Interval.shifted_pairing,MonadFormal.Temporal.Variational.Interval.unit_integral_eq_interval,MonadFormal.Temporal.Variational.Interval.unit_integral_indicator,MonadFormal.Temporal.Variational.Interval.shiftForce_pairing}
\end{proof}

% Lean module: MonadFormal.Temporal.Variational.Regularity
\subsubsection{Regularity from the weak equation}
\label{sec:variational-regularity}

\begin{lemma}[Mean-zero orthogonality and integration by parts]
\label[lemma]{lem:variational-orthogonality}
If $k\in V$ is orthogonal to every $h\in H$, then $k$ is
constant almost everywhere. For $g\in V$ and $h\in H$,
\[
 \langle Pg,h\rangle_{L^2(\mu)}
 =-\int g(t)(Ph)(t)\,d\mu(t).
\]
Here the bounded continuous function $Pg$ is also regarded as an
element of $V$.
\end{lemma}
\begin{proof}
Put $m=\int k\,d\mu$. Since $\mu([0,1])=1$, the function
$h=k-m$ belongs to $H$. Constants are orthogonal to $H$, so
$\lVert h\rVert_2^2=\langle k-m,h\rangle=0$ and $k=m$
almost everywhere. For the second identity, $Pg$ and $Ph$ are
absolutely continuous, with derivatives $g$ and $h$ almost
everywhere. Integration by parts has zero boundary term because
$Ph(0)=Ph(1)=0$.
\leandef{MonadFormal.Temporal.Variational.Interval.primitiveLp}
\leanproof{MonadFormal.Temporal.Variational.Interval.endpoint_const,MonadFormal.Temporal.Variational.Interval.const_inner_variation,MonadFormal.Temporal.Variational.Interval.eq_const_of_orthogonal,MonadFormal.Temporal.Variational.Interval.primitiveLp_ae,MonadFormal.Temporal.Variational.Interval.primitive_pairing}
\end{proof}

\begin{theorem}[Weak equation determines an absolutely continuous velocity]
\label{thm:variational-regularity}
Let $f\in H$, $g\in V$ and $\lambda\in\mathbb R$ satisfy
\[
 \langle f,h\rangle+\lambda\int g(t)(Ph)(t)\,d\mu(t)=0
 \qquad(h\in H).
\]
There is a constant $c$ such that $f=c+\lambda Pg$ almost
everywhere. Consequently, $x_f$ has a continuous velocity
\[
 v(t)=b-a+c+\lambda(Pg)(t),
\]
which is absolutely continuous on $[0,1]$. The equality
$x_f'(t)=v(t)$ holds at every interior point, and
$v'=x_f''=\lambda g$ almost everywhere on $[0,1]$.
\end{theorem}
\begin{proof}
By \Cref{lem:variational-orthogonality}, the weak equation says
that $f-\lambda Pg$ is orthogonal to $H$, hence is constant.
The primitive $Pg$ is continuous and absolutely continuous, so
the displayed formula gives the same regularity for $v$.
Integrating the continuous representative $c+\lambda Pg$ of $f$
upgrades the derivative of $Pf$ to every interior point by the
fundamental theorem of calculus. Finally $(Pg)'=g$ almost
everywhere. Since $x_f'=v$ on the open interval, their derivatives
also agree wherever $v'$ exists there.
\leanproof{MonadFormal.Temporal.Variational.Interval.weak_velocity_representation,MonadFormal.Temporal.Variational.Interval.primitive_hasDerivAt_of_continuous,MonadFormal.Temporal.Variational.Interval.weak_path_regularity,MonadFormal.Temporal.Variational.Interval.second_derivative_ae}
\end{proof}

% Lean module: MonadFormal.Temporal.Variational.Euler
\subsubsection{The nonlocal Euler equation}
\label{sec:variational-euler}

\begin{theorem}[Nonlocal Euler equation for the minimizing path]
\label{thm:nonlocal-euler}
Let $0\le\delta\le1$ and $\lambda\ge0$. The unique admissible
minimizer $x_*$ has a continuous interior velocity with an
absolutely continuous extension to $[0,1]$, and $x_*''\in L^2(\mu)$.
It satisfies, almost everywhere on $[0,1]$,
\begin{equation}\label{eq:nonlocal-euler}
 -x_*''(t)+\lambda G_\delta[x_*](t)=0.
\end{equation}
Where both shifted neighbours are admitted, the force is
\[
 G_\delta[x_*](t)=2x_*(t)-x_*(t-\delta)-x_*(t+\delta)
 \qquad(\delta\le t\le1-\delta).
\]
\end{theorem}
\begin{proof}
Take the minimizing coordinate from \Cref{thm:interval-minimum}.
By \Cref{lem:shift-pairing}, its weak equation is
\[
 \langle f_*,h\rangle+\lambda\int
 G_\delta[x_*](t)(Ph)(t)\,d\mu(t)=0\qquad(h\in H).
\]
The force belongs to $L^2(\mu)$. Apply
\Cref{thm:variational-regularity} with $g=G_\delta[x_*]$.
This gives the velocity regularity and
$x_*''=\lambda G_\delta[x_*]$ almost everywhere, including
square-integrability. The coordinate/path bijection identifies
this path with the unique minimizer in $\mathcal A_{a,b}$.
When $\delta\le t\le1-\delta$, both indicators in
\eqref{eq:variational-force} are one; expanding gives the last formula.
\leanproof{MonadFormal.Temporal.Variational.Interval.path_continuous,MonadFormal.Temporal.Variational.Interval.shift_inner_eq_force,MonadFormal.Temporal.Variational.Interval.exists_minimizer_euler,MonadFormal.Temporal.Variational.Interval.shiftForce_both}
\end{proof}

This is a strong equation almost everywhere, not a claim of two-sided
endpoint derivatives or a pointwise second derivative at every cutoff.
If $\delta>1/2$, the region with both neighbours is empty; the
indicator formula still applies. The delayed and advanced terms arise
from one boundary-constrained minimization problem. They do not supply
an initial-value rule that reads an independently chosen future.



